\documentclass[11pt]{article}
\usepackage[a4paper,margin=25mm]{geometry}
\usepackage[T1]{fontenc}
\usepackage{lmodern,amsmath,amssymb,amsthm,mathtools,booktabs,microtype}
\usepackage{enumitem}
\usepackage[noend]{algpseudocode}
\usepackage[dvipsnames]{xcolor}
\usepackage[colorlinks=true,linkcolor=MidnightBlue,urlcolor=MidnightBlue,citecolor=MidnightBlue]{hyperref}
\hypersetup{pdftitle={Extracting a crossing in time O(nm)},pdfauthor={MohammadTaghi Hajiaghayi and Mahdi JafariRaviz}}
\setlength{\parskip}{4pt}
\setlength{\parindent}{0pt}
\setlength{\emergencystretch}{2em}

\newtheorem*{theorem}{Theorem}
\newtheorem{lemma}{Lemma}
\newtheorem*{factT}{Fact T}
\newtheorem*{factP}{Fact P}

\DeclareMathOperator{\Planar}{Planar}
\DeclareMathOperator{\BridgeCross}{BridgeCross}
\DeclareMathOperator{\Classify}{Classify}
\newcommand{\cR}{\mathcal R}
\newcommand{\vin}[1]{#1_{\mathrm{in}}}
\newcommand{\vout}[1]{#1_{\mathrm{out}}}
\newcommand{\ans}[1]{\texttt{#1}}

\title{\textbf{Extracting a crossing in time $O(nm)$}}
\author{MohammadTaghi Hajiaghayi \and Mahdi JafariRaviz}
\date{November 2026}

\begin{document}
\maketitle

\begin{abstract}
Let $C$ be a cycle in a finite simple graph $G$ with $n$ vertices and $m$
edges. We give an algorithm that either finds two vertex-disjoint $C$-paths
whose ends alternate around $C$, or correctly reports that there are none. It
runs in time $O(nm)$. The algorithm deletes parts of the graph behind
separators of order at most three, contracts edges, and lifts the crossing it
finds in the small graph back to $G$. Besides graph searches and bounded flows,
it uses one more primitive: a Boolean linear-time planarity test. In the Lean
formalization of the Graph Minors series, this algorithm replaces the two
paths algorithm of Shiloach that Graph Minors XIII uses in (8.1).
\end{abstract}

\section*{Introduction}

All graphs in this note are finite and simple. For a set $U$ of vertices of a
graph $K$, we write $N_K(U)$ for the set of vertices outside $U$ that have a
neighbor in $U$. For a set $X$ of vertices, $K[X]$ is the subgraph induced by
$X$.

Let $C$ be a cycle of a graph $G$. A \emph{$C$-path} is a path with at least
one edge whose two ends lie on $C$ and whose other vertices do not. A
\emph{crossing} of $C$ is a pair of vertex-disjoint $C$-paths, one with ends
$a,b$ and the other with ends $c,d$, such that $a,c,b,d$ occur in this cyclic
order around $C$. The cycle $C$ is \emph{crossed} if it has a crossing, and
\emph{crossless} otherwise. These are the definitions of Humeau and Pous
\cite[Definition~2.4, page~7]{HP25}.

\begin{theorem}
There is an algorithm that, given a graph $G$ with $n$ vertices and $m$ edges
and a cycle $C$ of $G$, returns a crossing of $C$ if $C$ is crossed, and
reports that $C$ is crossless otherwise. Its running time is $O(nm)$.
\end{theorem}

\paragraph{Where the theorem is used.}
Algorithm (8.1) of Graph Minors XIII \cite[page~93]{GM13} finds a cross in a
society. It builds an auxiliary graph and applies an algorithm of Shiloach
\cite{Sh80} to find two vertex-disjoint paths joining $s_1,t_1$ and $s_2,t_2$.
Shiloach's paper describes the algorithm through a sequence of reductions, and
on page~454 it says: ``We have not written the solution as an explicit, long,
and tedious algorithm.'' For distinct vertices $s_1,s_2,t_1,t_2$ of a graph
$H$, two vertex-disjoint paths joining $s_1,t_1$ and $s_2,t_2$ are exactly a
crossing of the cycle $s_1s_2t_1t_2$ in the graph obtained from $H$ by adding
the four edges of this cycle \cite[Section~1.2, pages~2--3]{HP25}. So the
theorem supplies this step of (8.1). In the Lean formalization of the Graph
Minors series, the algorithm of this note is used in place of Shiloach's
algorithm (Section~\ref{sec:lean}).

\paragraph{Outline.}
The algorithm runs once on $(G,C)$. Section~\ref{sec:facts} states the facts we
use from outside. Sections~\ref{sec:root} to~\ref{sec:clean} give a flow
primitive, a reversible replacement at a separator of order at most three, and
an initial cleaning. Sections~\ref{sec:bridges} and~\ref{sec:classify} classify
a clean instance in linear time. The classification either finds a crossing,
or shows that there is none, or finds that the graph is nonplanar, which
implies that the cycle is crossed. In the last case, Sections~\ref{sec:sparse}
to~\ref{sec:loop} shrink the graph by contracting edges, until the
classification finds a crossing. The main technical point is that a trial
contraction can be checked and repaired with a bounded number of flow searches
(Sections~\ref{sec:twocut} and~\ref{sec:repair}). Section~\ref{sec:lift} lifts
the crossing back to $G$, and Section~\ref{sec:time} counts the work.
Section~\ref{sec:remarks} lists the assumptions and describes the reference
implementation \texttt{crossing\_algorithm.py} and its tests.
Section~\ref{sec:lean} describes the Lean formalization.

\section{The external facts used}\label{sec:facts}

Besides Menger's theorem for vertex-disjoint paths and the elementary flow
computation of Section~\ref{sec:flow}, we use the following facts.

\subsection*{The two paths theorem: the consequence we need}

\begin{factT}
Let $C$ be a cycle of a graph $K$. If $C$ is crossless and $K$ is nonplanar,
then there is a nonempty set $U\subseteq V(K)\setminus V(C)$ with
$|N_K(U)|\le 3$.
\end{factT}

Fact~T follows from the generalized two paths theorem of Humeau and Pous
\cite[Theorem~1.3, page~2]{HP25}. To state it, we recall webs
\cite[Definition~2.1, page~6]{HP25}. A \emph{rib} with frame $C$ is a planar
graph with a plane embedding in which the outer face is bounded by the cycle
$C$, every inner face is a triangle, and every triangle bounds an inner face.
A \emph{web} with frame $C$ is obtained from a rib with frame $C$ by adding,
for each triangle $T$ of the rib, a clique, possibly with no vertices, whose
vertices are adjacent to the three vertices of $T$ and to no other vertex
outside the clique. A graph $H$ with a cycle $C$ is \emph{maximally
$C$-crossless} if $C$ is crossless in $H$ but crossed in every graph obtained
from $H$ by adding one edge. The generalized two paths theorem says that $H$ is
a web with frame $C$ if and only if $H$ is maximally $C$-crossless. We use only
the direction proved in \cite[Lemma~4.3, page~13]{HP25}: a maximally
$C$-crossless graph is a web with frame $C$.

\begin{proof}[Proof of Fact T]
Suppose that $C$ is crossless in $K$. Add edges between vertices of $K$, one at
a time, as long as $C$ stays crossless. Since $K$ is finite, this stops with a
maximally $C$-crossless graph $H$ on the vertex set of $K$, and $H$ is a web
with frame $C$. If no clique of the web has a vertex, then $H$ is its rib,
which is planar, and so is its subgraph $K$. Since $K$ is nonplanar, some
clique, filling a triangle $T$, has a vertex. Let $U$ be the vertex set of
this clique. Then $U$ is nonempty, and it is disjoint from $V(C)$ because $C$
consists of rib vertices. Since $K$ is a subgraph of $H$ on the same vertex
set, $N_K(U)\subseteq N_H(U)=V(T)$, and so $|N_K(U)|\le3$.
\end{proof}

This is the only consequence of the two paths theorem that we use.

\subsection*{The Boolean planarity primitive}

\begin{factP}
On an adjacency-list representation of a graph $K$, a deterministic procedure
$\Planar(K)$ returns whether $K$ is planar in time $O(|V(K)|+|E(K)|)$.
\end{factP}

The procedure need not return an embedding or a Kuratowski subdivision.
Hopcroft and Tarjan \cite{HT74} gave the first linear-time planarity test. The
companion note \cite{HJ26} proves a Boolean left-right test of this kind in
full, and the reference implementation uses that test.

\subsection*{Euler's bound}

We also use Euler's bound: a planar graph with $N\ge3$ vertices has at most
$3N-6$ edges, and a bipartite planar graph with $N\ge3$ vertices has at most
$2N-4$ edges. In particular, $K_{3,3}$ is nonplanar: a bipartite planar graph
on six vertices has at most eight edges, whereas $K_{3,3}$ has nine.

\section{The root condition and the bounded-flow primitive}\label{sec:root}

An \emph{instance} $(K,D)$ consists of a graph $K$ and a cycle $D$ of $K$. Add
a new vertex $x$, adjacent to every vertex of $D$, and write $J=K+x$. We call
$x$ the \emph{apex}. It is used only in auxiliary computations.

For $k\ge1$, we say that $\cR_k(J,x)$ holds if $J-S$ is connected for every set
$S\subseteq V(J)\setminus\{x\}$ with $|S|<k$. All these separators exclude
$x$. We do not claim ordinary $k$-connectivity, which would also consider
separators that contain $x$.

We call $(K,D)$ \emph{clean} when $\cR_4(K+x,x)$ holds. A \emph{four-fan} from
a vertex $v$ to $D$ consists of four paths that share only their first vertex
$v$ and end at four distinct vertices of $D$. Because every vertex of $D$ is
adjacent to $x$, Menger's theorem gives
\begin{equation*}
(K,D)\text{ is clean}
\quad\Longleftrightarrow\quad
\text{every vertex }v\notin V(D)\text{ has a four-fan to }D.
\end{equation*}
Indeed, $\cR_4(J,x)$ says that no vertex is separated from $x$ by at most three
vertices other than itself and $x$. A vertex of $D$ is adjacent to $x$. For a
vertex $v\notin V(D)$, which is not adjacent to $x$, Menger's theorem turns
this into four paths from $v$ to $x$ that share only $v$ and $x$. Cutting each
of these paths at its first vertex of $D$ gives a four-fan, and extending the
paths of a four-fan by their edges to $x$ gives four such paths back.

\subsection{Exact bounded-flow computation}\label{sec:flow}

Every flow query below has a source $z$ and the sink $x$, where $z$ is not
adjacent to $x$, and asks for $k\le4$ internally vertex-disjoint paths from
$z$ to $x$. We build the split network.
\begin{itemize}[nosep]
\item Replace each vertex $v$ by two nodes $\vin v$ and $\vout v$.
\item Give the arc $\vin v\to\vout v$ capacity $1$, except at $z$ and $x$,
where its capacity is $k+1$.
\item For every edge $uv$, add the arcs $\vout u\to\vin v$ and $\vout v\to\vin u$,
each of capacity $k+1$.
\item Use the source $\vout z$ and the sink $\vin x$.
\end{itemize}
Starting with the zero flow, we search the residual graph for a path from the
source to the sink and augment one unit along it. We stop after $k$
augmentations or after an unsuccessful search. Thus each query uses at most
$k$ residual searches, each on a graph of size $O(|V|+|E|)$. Any augmenting
paths work.

If the computation stops below $k$, its minimum cut consists of arcs
$\vin v\to\vout v$ of capacity $1$, because every other arc has capacity $k+1$,
more than the value of the cut. These arcs give a vertex separator of that size
between $z$ and $x$.

\subsection{Computing a minimum cut with maximum source side}\label{sec:maxcut}

One later step, Lemma~\ref{lem:twocut}, needs a particular minimum separator.

After a failed bounded-flow computation, let $W$ be the set of nodes of the
network that \emph{cannot reach the sink in the residual graph}. We compute
$W$ by a search from the sink along reversed residual arcs, and take the
complement. By a standard property of maximum flows, $W$ is the largest source
side of a minimum network cut: it contains the source side of every minimum
cut. Let
\begin{equation*}
S=\{v:\ \vin v\in W,\ \vout v\notin W\},
\end{equation*}
let $R_x$ be the component of $x$ in $J-S$, and let
\begin{equation*}
U=V(J)\setminus\bigl(S\cup R_x\bigr).
\end{equation*}
In general, the \emph{open source side} of a vertex separator $S'$ between $z$
and $x$ is the set of vertices outside $S'$ that are not in the component of
$x$ in $J-S'$. So $U$ is the open source side of $S$.

\emph{This $U$ contains the open source side of every minimum vertex separator
between $z$ and $x$.} To see this, let $S'$ be a minimum vertex separator
between $z$ and $x$, with open source side $U'$. The nodes $\vin u,\vout u$ for
$u\in U'$, together with the nodes $\vin s$ for $s\in S'$, form the source side
of a network cut of capacity $|S'|$. So this cut is minimum, and its source
side is contained in $W$. Now take a path of $J$ from a vertex of $U'$ to $x$
that avoids $S$. The corresponding path in the network starts in $W$ and ends
at the sink, which is not in $W$, so it uses an arc that leaves $W$. That arc is
not an arc $\vin v\to\vout v$ with $v\notin S$, by the definition of $S$. It is
not an arc of capacity $k+1$ either, since such an arc would make the capacity
of the minimum cut larger than its value. This is a contradiction. Hence no
vertex of $U'$ is joined to $x$ in $J-S$, and no vertex of $U'$ is in $S$,
because its node $\vout u$ lies in $W$. Therefore $U'\subseteq U$.

This fixes the choice of the cut completely. \emph{An arbitrary minimum cut
does not suffice for Lemma~\ref{lem:twocut}.}

\section{A reversible replacement at a small separator}\label{sec:replace}

The following operation is used throughout.

Let $A\subseteq V(K)\setminus V(D)$ be a connected set, and let $S=N_K(A)$,
with $|S|\le 3$. The \emph{replacement} deletes $A$ and adds every missing edge
of the clique on $S$. We record the graph before the operation, the set $A$,
and the set $S$.

\begin{lemma}[Replacement preserves crossing existence]\label{lem:replace}
If $D$ is crossed before the replacement, then it is crossed afterward. A
crossing afterward can be lifted to a crossing before the replacement with at
most one ordinary path search.
\end{lemma}

\begin{proof}
\emph{Forward direction.} A crossing path that enters $A$ must use at least two
distinct vertices of $S$, since neither of its ends lies in $A$. Since
$|S|\le3$, the two disjoint crossing paths cannot both enter $A$. For the path
that enters $A$, if there is one, replace the segment between its first and
last visits to $S$ by the clique edge between these two vertices. This
introduces no new vertices, keeps the ends of the path, and keeps the two paths
disjoint.

\emph{Reverse direction.} If neither returned path uses a newly added edge, we
keep both. Otherwise exactly one path uses new edges, because two disjoint
paths that both use new clique edges would need four distinct vertices of $S$.
Let $p$ and $q$ be the first and last vertices of $S$ on this path. Here
$p\ne q$, since a new edge joins two vertices of $S$. We search for a path from
$p$ to $q$ in $K[A\cup\{p,q\}]$. Such a path exists because $A$ is connected
and both $p$ and $q$ have neighbors in $A$. We replace the whole segment
between $p$ and $q$ by this path. Its inner vertices lie in $A$, which neither
returned path used, and they include no vertex of $S$ or $D$. Thus the new
path is simple, the two paths stay disjoint, the ends do not change, and they
still alternate around $D$.
\end{proof}

We also allow deleting several components of $K-S$ at once, all with
neighborhood $S$, provided we record one of them. The recorded component is
connected and has a neighbor at every vertex of $S$. The forward proof is
unchanged, and the reverse search uses the recorded component.

\begin{lemma}[Replacement preserves connectivity statements on retained vertices]\label{lem:retain}
Let $T$ be a set of retained vertices. Every path of the old graph that avoids
$T$ and has retained ends gives a path of the new graph between the same ends
that avoids $T$. This holds with or without the apex.
\end{lemma}

\begin{proof}
Replace each excursion of the path through a deleted component by the edge of
the boundary clique between the vertices where the path enters and leaves the
component, or delete the excursion if these two vertices are equal. These
boundary vertices lie on the path, so they are not in $T$. The result is a walk
of the new graph that avoids $T$, and shortcutting repeated vertices gives a
path. The apex is not affected, because a deleted component contains no vertex
of $D$.
\end{proof}

Therefore every statement proved earlier that a retained vertex cannot be
separated from $x$ by at most a given number of other vertices stays true
after the replacement. Lemma~\ref{lem:retain} concerns ordinary connectivity
and does not use the two paths theorem.

\section{Initial cleaning}\label{sec:clean}

First we discard the components of $G$ that do not contain $C$, since they
cannot contain a $C$-path. We start with the instance $(K,D)$, where $K$ is the
component of $G$ that contains $C$, and $D=C$.

Then we process the vertices of $K$ that are not on $C$ in any fixed order,
skipping vertices that have already been deleted. For each such vertex $v$
that is still present:
\begin{enumerate}[nosep]
\item In $K+x$, run the bounded-flow computation asking for four paths from
$v$ to $x$.
\item If it succeeds, continue with the next vertex.
\item Otherwise, obtain a separator $S$ of size at most three.
\item Find the component $A$ of $K-S$ that contains $v$.
\item Replace $A$ by the clique on its actual neighborhood $N_K(A)$, using
Lemma~\ref{lem:replace}.
\end{enumerate}

The component $A$ contains no cycle vertex. Indeed, every cycle vertex that is
still present is adjacent to $x$, whereas $S$ separates $v$ from $x$. Since
$N_K(A)\subseteq S$, Lemma~\ref{lem:replace} applies.

Every successful four-path query stays true afterward, by
Lemma~\ref{lem:retain}. Every unsuccessful query deletes the vertex it asked
about. Thus, after at most $n$ bounded-flow queries, every remaining vertex off
the cycle has four paths to $x$ that share only their ends, and the remaining
instance is clean by the equivalence of Section~\ref{sec:root}.

Every replacement deletes at least one vertex and adds at most three edges. We
record the replacements for the lifting at the end.

\section{A linear procedure for crossings in distinct cycle bridges}\label{sec:bridges}

We need the following procedure for any instance $(K,D)$:
\begin{quote}
Return a crossing whose two paths use distinct $D$-bridges, or correctly report
that no such crossing exists.
\end{quote}
A \emph{$D$-bridge} is either
\begin{itemize}[nosep]
\item a component of $K-V(D)$, together with its edges to $D$, or
\item a single \emph{chord}: an edge that is not an edge of $D$ and has both
ends on $D$.
\end{itemize}
Different component bridges have disjoint sets of inner vertices, and a chord
is its own bridge. Every $D$-path is a chord or has all its inner vertices in
one component of $K-V(D)$, so it uses exactly one bridge. The
\emph{attachments} of a component bridge are the vertices of $D$ adjacent to
the component, and the attachments of a chord are its two ends. The ends of a
$D$-path are attachments of the bridge it uses.

\subsection{Candidate endpoint pairs}\label{sec:candidates}

We compute the attachment list of every bridge in cyclic order. This takes time
$O(n+m)$: compute the components off the cycle, then scan the vertices of $D$
in order and append the identifiers of their incident components, suppressing
repeated incidences at the same cycle vertex.

For a bridge with attachment list $t_0,t_1,\ldots,t_{h-1}$, we generate the
pairs
\begin{equation*}
\{t_i,t_{i+1}\},
\qquad
\{t_i,t_{i+2}\},
\end{equation*}
with indices modulo $h$, and omit pairs whose two elements are equal.
Duplicates may be kept. These pairs are the \emph{candidates} of the bridge.
There are at most twice as many candidates as attachments.

\begin{lemma}[These candidates detect every crossing between two bridges]\label{lem:candidates}
Let $A$ and $B$ be the attachment sets of two distinct bridges. Suppose that
$A$ contains $a_1,a_2$ and $B$ contains $b_1,b_2$ such that the pairs
$\{a_1,a_2\}$ and $\{b_1,b_2\}$ alternate around $D$. Then some candidate of
the first bridge alternates with some candidate of the second.
\end{lemma}

\begin{proof}
There are $A$-attachments in both open arcs of $D$ determined by $b_1$ and
$b_2$. Around $b_2$, choose the nearest $A$-attachment strictly before it and
the nearest one strictly after it. These two attachments lie in the two
different open arcs, so they alternate with $b_1,b_2$. In the cyclic
$A$-attachment list they are consecutive, unless $b_2\in A$, in which case
exactly one attachment, namely $b_2$, lies between them. Hence they form an
$A$-candidate that alternates with $b_1,b_2$.

Applying the same argument to $B$, relative to this $A$-candidate, gives a
$B$-candidate that alternates with it.
\end{proof}

The candidates at distance two are needed to handle shared attachments.

\subsection{Finding differently colored alternating candidates in linear time}\label{sec:sweep}

Number the vertices of $D$ as $0,\ldots,c-1$ in cyclic order. We represent a
candidate as a triple
\begin{equation*}
(l,r,\gamma),\qquad l<r,
\end{equation*}
where $l$ and $r$ are the numbers of its two attachments and $\gamma$, its
\emph{color}, is the identifier of its bridge. Two candidates $(l',r',\gamma')$
and $(l,r,\gamma)$ alternate around $D$ exactly when $l'<l<r'<r$ or
$l<l'<r<r'$.

Sort the candidates by increasing $l$, breaking ties by decreasing $r$. Two
stable counting sorts do this in linear time.

We maintain a stack of blocks, each block containing candidates of one color,
with the following invariants.
\begin{enumerate}[nosep]
\item Adjacent blocks have different colors.
\item Every right endpoint in a block is at most every right endpoint in the
block immediately below it.
\item Before a candidate with left endpoint $l$ is processed, every stored
candidate has right endpoint greater than $l$.
\end{enumerate}

For a new candidate $I=(l,r,\gamma)$:
\begin{enumerate}[nosep]
\item Delete the candidates with right endpoint at most $l$ from the top block.
Pop the top block if it becomes empty, and repeat with the new top block.
\item Let $B$ be the first block from the top whose color differs from
$\gamma$. By invariant~1, this is the top block or the second block.
\item If $B$ exists and its minimum right endpoint $r'$ satisfies $r'<r$,
return $I$ and the candidate in $B$ that attains this minimum.
\item Otherwise, append $I$ to the top block if its color is $\gamma$, and
push a new block containing $I$ if not.
\end{enumerate}

\paragraph{Correctness.}
In step~3, the returned earlier candidate $(l',r',\gamma')$ satisfies
\begin{equation*}
l'<l<r'<r.
\end{equation*}
Here $l'\le l$ by the sorting order, $l<r'$ by invariant~3, and $r'<r$ by the
test. The strict inequality $l'<l$ follows from the tie-breaking order: an
earlier candidate with the same left endpoint would have right endpoint at
least $r$. So the two candidates alternate, and they have different colors.

The endpoint-order invariant~2 makes the tested minimum no greater than the
right endpoint of any stored candidate of a different color. Thus a failed
test excludes every differently colored crossing between the new candidate and
a stored one. A candidate that alternates with a later one has right endpoint
greater than the left endpoint of the later one, so step~1 has not deleted it
when the later one is processed. Therefore the sweep finds a differently
colored alternating pair whenever one exists. A failed test also gives exactly
the inequality needed to preserve invariant~2 when appending or pushing.

Cleanup needs to inspect only the top block. If that block keeps an active
candidate, then by invariant~2 every right endpoint in every lower block is at
least its right endpoint, so these candidates are active too.

\paragraph{Linear bound.}
At any sweep position, one bridge has at most six active candidates. A cut of
the cyclic attachment list at the sweep position separates it into two
consecutive blocks. At most two pairs at cyclic distance one and four pairs at
cyclic distance two cross this cut, counting duplicates. Therefore every stack
block has size at most six. Each candidate is inserted once and removed at most
once.

\paragraph{The two paths.}
When an alternating pair of candidates of different colors is found, we make
two searches. For a component bridge with component $A$ and selected endpoints
$p,q$, we search in $K[A\cup\{p,q\}]$. For a chord, we use its edge. The two
searches may return arbitrary paths: their inner vertices belong to different
bridges, and the four selected endpoints are distinct.

We call this whole procedure $\BridgeCross(K,D)$. It takes time $O(n+m)$ and,
when it succeeds, two path searches.

\subsection{Why this procedure is complete on planar graphs}\label{sec:lemma4}

Suppose that a crossing $Q,R$ has the inner vertices of both paths in the same
component $A$ of $K-V(D)$. Connect an inner vertex of $Q$ to an inner vertex of
$R$ by a path inside $A$, and shorten it so that its inner vertices avoid
$Q\cup R$.

The cycle, the two crossing paths and this connector form a subdivision of
$K_{3,3}$. Let $q$ and $r$ be the ends of the connector on $Q$ and $R$, and let
$q_1,r_1,q_2,r_2$ be the ends of the crossing in cyclic order, with $q_1,q_2$
the ends of $Q$. The bipartition of the six branch vertices is
\begin{equation*}
\{q_1,q_2,r\},
\qquad
\{r_1,r_2,q\}.
\end{equation*}
The nine branch paths are the four arcs of the cycle, the four halves of $Q$
and $R$, and the connector.

This is impossible in a planar graph, by Euler's bound in
Section~\ref{sec:facts}. A chord already forms a separate bridge.
Consequently:

\begin{lemma}\label{lem:planar}
Every crossing in a planar graph uses distinct cycle bridges. Therefore
$\BridgeCross$ is complete on planar inputs.
\end{lemma}

\section{Classifying a clean instance}\label{sec:classify}

For a clean instance $(K,D)$, we run the following procedure.
\begin{algorithmic}[1]
\Function{Classify}{$K,D$}
\If{$\BridgeCross(K,D)$ returns a crossing}
\State \Return \ans{WITNESS} and that crossing
\EndIf
\If{$\Planar(K)$}
\State \Return \ans{NO}
\EndIf
\State \Return \ans{YES}
\EndFunction
\end{algorithmic}

The answer \ans{YES} means that a crossing exists, but no crossing has been
found yet. The classification is exact.
\begin{itemize}[nosep]
\item \ans{WITNESS} is correct by construction.
\item \ans{NO} is correct by Lemma~\ref{lem:planar}.
\item For \ans{YES}, the graph $K$ is nonplanar. If $D$ were crossless, Fact~T
would give a nonempty set $U$ of vertices off the cycle with
$|N_K(U)|\le3$. Then $N_K(U)$ would separate $U$ from $x$ in $K+x$,
contradicting cleanliness.
\end{itemize}

Thus classification takes linear time. The remaining work is to turn a
nonplanar \ans{YES} answer into paths by making the graph smaller.

\section{Sparsifying without losing a positive answer}\label{sec:sparse}

We need a sparse graph, so that a bounded number of vertices of bounded degree
can be examined.

\subsection{Four BFS forests preserve all cuts of order at most three}\label{sec:forests}

A breadth-first search scans vertices in the order of a queue. When it scans a
vertex $s$, it marks every unmarked neighbor of $s$, selects the edge joining
it to $s$, and appends it to the queue. A \emph{BFS spanning forest} of a graph
consists of the selected edges, where a new search, with a new root, starts at
each vertex that is still unmarked when its turn comes in a fixed order.

Given a graph $J$, we compute four forests in succession.
\begin{enumerate}[nosep]
\item Take a BFS spanning forest $F_1$ of $J$.
\item Delete its edges and take a BFS spanning forest $F_2$ of the rest.
\item Delete those edges and take $F_3$.
\item Delete those edges and take $F_4$.
\end{enumerate}
Let $F=F_1\cup F_2\cup F_3\cup F_4$.

This is the scan-first construction of sparse certificates of Cheriyan, Kao and
Thurimella \cite{CKT93}. The following proof gives the exact strong property
that we need.

\begin{lemma}[Strong certificate]\label{lem:forests}
Let $F$ be the union of $k$ successive BFS forests of $J$, taken as above. For
every set $S$ of vertices with $|S|<k$, the graphs $F-S$ and $J-S$ have
exactly the same components.
\end{lemma}

\begin{proof}
Suppose otherwise. Then there is a partition
\begin{equation*}
V(J)=A\;\dot\cup\;S\;\dot\cup\;B
\end{equation*}
such that $F$ has no edge between $A$ and $B$, but $J$ has an edge $ab$ with
$a\in A$ and $b\in B$. Take $A$ to be one component of $F-S$ that is not a
component of $J-S$, and let $B$ contain all other components.

Before the forest $F_i$ is computed, we maintain parts $A_i\supseteq A$ and
$B_i\supseteq B$ and a remaining separator $S_i$, which partition $V(J)$, with
the following properties:
\begin{itemize}[nosep]
\item the vertex $a$ is in $A_i$, and $b$ is in $B_i$;
\item the edge $ab$ is still in the residual graph, that is, it has not been
selected by an earlier forest;
\item no edge of the current or any later forest joins $A_i$ to $B_i$.
\end{itemize}
Initially $A_1=A$, $B_1=B$ and $S_1=S$.

In the residual component that contains $ab$, the current spanning tree
connects $a$ to $b$. Since it has no edge between $A_i$ and $B_i$, it contains
a vertex of $S_i$. Let $s$ be the first such vertex scanned.

Suppose that the root of this tree lies in $A_i$. Then no vertex of $B_i$ was
marked before $s$ was scanned. Otherwise, the tree path from the root to that
vertex would contain a proper ancestor in $S_i$, and this ancestor was scanned
before that vertex was marked, contradicting the choice of $s$.

Therefore, when $s$ is scanned, every residual neighbor of $s$ in $B_i$ is
unmarked, and the search selects every edge from $s$ to $B_i$. We move $s$ into
$A_i$. All crossing edges that this move creates have just been selected and
removed, so no later forest contains them.

The case where the root lies in $B_i$ is symmetric. If the root lies in $S_i$,
it is $s$. All its neighbors are unmarked when it is scanned, so all its
incident edges are selected, and $s$ can be moved to either side.

Each forest removes one vertex from the separator and preserves the properties. The edge
$ab$ is never selected, because it always joins the two parts. After $|S|$
forests the separator is empty. Since $|S|<k$, there is a next forest. It must
connect $a$ and $b$, but it has no edge between the two parts, which now cover
$V(J)$. This is a contradiction.

The proof uses simplicity when $s$ is scanned: marking an unmarked neighbor
selects its unique edge from $s$.
\end{proof}

The four passes take time $O(n+m)$. On $n+1$ vertices, their union has at most
$4n$ edges.

\subsection{The sparsification step}\label{sec:sparsify}

Suppose that the current clean graph $K$ is nonplanar and has $N$ vertices and
$M$ edges.

If $M\le16N$, we do nothing. Otherwise:
\begin{enumerate}[nosep]
\item Compute the four-forest certificate of $J=K+x$.
\item Keep its edges that do not contain the apex.
\item Keep all edges of $D$.
\item Keep any $3N-5$ edges of $K$.
\item Delete all other edges of $K$.
\end{enumerate}

The kept graph has at most
\begin{equation*}
4N+N+(3N-5)=8N-5
\end{equation*}
edges.

Adding all edges from the apex to the cycle gives a graph that contains the
certificate and is contained in $J$. By Lemma~\ref{lem:forests}, removing at
most three vertices other than $x$ leaves the same components in all three
graphs, so the kept graph is still clean. Its $3N-5$ kept edges of $K$ make it
nonplanar by Euler's bound. As in Section~\ref{sec:classify}, Fact~T then
shows that $D$ is still crossed.

This step only deletes edges, so a crossing found later needs no special
lifting across it.

\section{The maximum two-cut lemma}\label{sec:twocut}

This is the central separator lemma.

\begin{lemma}\label{lem:twocut}
Let $B$ be a graph with a vertex $x$ such that $\cR_2(B,x)$ holds. Let $z\ne x$
be a vertex whose minimum vertex separators from $x$ have size two. Choose
$S=\{p,q\}$ with maximum open source side
\begin{equation*}
U=V(B)\setminus\bigl(S\cup R_x\bigr),
\end{equation*}
as in Section~\ref{sec:maxcut}, where $R_x$ is the component of $x$ in $B-S$.
Then neither $p$ nor $q$ can be separated from $x$ by at most two other
vertices.
\end{lemma}

\begin{proof}
Suppose that $p$ can be separated from $x$ by at most two other vertices.
Since $\cR_2(B,x)$ holds, a minimum such separator $T$ has size exactly two.
Let $Y$ be the union of the components of $B-T$ that do not contain $x$, so
$p\in Y$.

Consider the separations
\begin{align*}
(A_1,B_1)&=\bigl(U\cup S,\ V(B)\setminus U\bigr),\\
(A_2,B_2)&=\bigl(Y\cup T,\ V(B)\setminus Y\bigr).
\end{align*}
Their join and meet are
\begin{equation*}
\bigl(A_1\cup A_2,\ B_1\cap B_2\bigr),
\qquad
\bigl(A_1\cap A_2,\ B_1\cup B_2\bigr).
\end{equation*}
They are separations. Their orders add up to four, since a vertex-by-vertex
count shows that the orders of the join and the meet add up to the orders of
the two given separations.

The meet has order at least two.
\begin{itemize}[nosep]
\item If $U\cap Y\ne\varnothing$, the open left side of the meet is nonempty,
and $x$ lies on its open right side. A meet of order zero or one would
contradict $\cR_2(B,x)$.
\item If $U\cap Y=\varnothing$, the minimality of $S$ gives a neighbor $h\in U$
of $p$, since otherwise $\{q\}$ would separate $z$ from $x$. Since $p\in Y$
and $h\notin Y$, the edge $ph$ forces $h\in T$. The two distinct vertices $p$
and $h$ both belong to the separator of the meet.
\end{itemize}

Therefore the join has order at most two. Its open source side is $U\cup Y$,
which contains $z$, does not contain $x$, and strictly contains $U$, because
$p\in Y\setminus U$.

So the separator of the join is a minimum separator between $z$ and $x$, and
its open source side contains $U\cup Y$, which strictly contains $U$. This
contradicts the choice of $U$ in Section~\ref{sec:maxcut}. The argument for $q$
is the same.
\end{proof}

The graph $B$ need not be $2$-connected: $x$ may be a cut vertex. Every
separator used in the proof excludes $x$.

\section{Repairing one trial contraction with boundedly many searches}\label{sec:repair}

Start with a clean instance $(K,D)$. Choose an edge $uv$ of one of these forms:
\begin{itemize}[nosep]
\item the vertex $u$ is not on $D$, and $v$ is arbitrary;
\item the vertices $u$ and $v$ are both on $D$, and they are consecutive on $D$.
\end{itemize}
Contract $uv$ to a new vertex $w$, suppressing loops and duplicate edges. If
$u\notin V(D)$, the cycle is unchanged. If both ends are on $D$, remove $u$
from the cyclic list and keep the name $v$ for $w$.

Let $H$ be the contracted graph and $J=H+x$. Set
\begin{equation*}
M=N_K(u)\setminus\{v\}.
\end{equation*}
These are the \emph{original markers}. They stay fixed for the whole trial.

\subsection{What the contraction guarantees}\label{sec:guarantees}

The old augmented graph $K+x$ satisfied $\cR_4$. Therefore:
\begin{enumerate}[nosep]
\item The property $\cR_3(J,x)$ holds. A separator of size at most two that contains $w$
lifts to at most three old vertices, namely $u$, $v$ and the rest of the
separator. A separator that avoids $w$ lifts unchanged.
\item Every separator of size at most three that excludes $x$ and disconnects
$J$ contains $w$.
\item For
\begin{equation*}
B=J-w,
\end{equation*}
the property $\cR_2(B,x)$ holds.
\end{enumerate}
In each case, a separator of $J$ that avoids $w$ also separates $K+x$, because
contracting the edge $uv$ does not change whether the rest of the graph is
connected. Property~3 follows from property~1, since a separator of $B$ of size
at most one, together with $w$, would separate $J$.

Furthermore, every component $Q$ of $J-\{w,p,q\}$ that does not contain $x$,
for vertices $p,q$, contains a marker. Indeed, if $Q$ had no old neighbor of
$u$, then in the old augmented graph all neighbors of $Q$ would lie in
$\{v,p,q\}$. That would violate the old $\cR_4$.

Each replacement made below changes $J$ and $B$. Properties 1, 2 and 3 say that
certain vertices are joined to $x$ after certain sets of retained vertices are
deleted. So, by Lemma~\ref{lem:retain}, they stay true for the current $J$ and
$B=J-w$ after every replacement.

\subsection{Exact restoration procedure}\label{sec:restore}

We process the fixed list $M$, once for each marker. For a marker $z$:
\begin{enumerate}[nosep]
\item If $z$ has been deleted, skip it.
\item If $z$ is on the current cycle, skip it. Its edge to $x$ keeps it out of
every component that does not contain $x$.
\item Form $B=J-w$.
\item Run the bounded-flow computation for three internally disjoint paths
from $z$ to $x$ in $B$.
\item If it succeeds, continue with the next marker.
\item Otherwise, obtain the minimum two-cut $\{p,q\}$ with maximum open source
side $U$, as in Section~\ref{sec:maxcut}.
\item Delete all vertices of $U$ from $H$, and complete the triangle
$\{w,p,q\}$.
\item Record one component of $B-\{p,q\}$ inside $U$, for example the
component containing $z$, as the connected representative for the later
lifting.
\end{enumerate}

The failed flow has value exactly two, by $\cR_2(B,x)$.

\subsection{Why each replacement is valid}\label{sec:valid}

Each deleted component $A$ of $B-\{p,q\}$ inside $U$ has exactly the
neighborhood
\begin{equation*}
N_J(A)=\{w,p,q\}.
\end{equation*}
It has a neighbor at both $p$ and $q$: otherwise one vertex would separate it
from $x$ in $B$, violating $\cR_2(B,x)$. It has a neighbor at $w$: otherwise
$\{p,q\}$ would separate it from $x$ in $J$, violating $\cR_3(J,x)$. It
contains no vertex of the current cycle, since every cycle vertex outside
$\{w,p,q\}$ is still adjacent to $x$.

Thus Lemma~\ref{lem:replace} applies, in its version that deletes several
components with a common neighborhood.

\subsection{Why no additional markers are needed}\label{sec:markers}

We call a retained vertex \emph{good} if no set of at most three other
vertices, excluding $x$, separates it from $x$. Otherwise it is \emph{bad}. A
\emph{bad component} is a component, not containing $x$, of the graph obtained
from $J$ by deleting at most three vertices other than $x$. All its vertices
are bad.

At every replacement made during restoration, the three boundary vertices
$w,p,q$ are good.
\begin{itemize}[nosep]
\item The vertex $w$ is good, because by property~2 of
Section~\ref{sec:guarantees} every separator of size at most three contains
$w$.
\item By Lemma~\ref{lem:twocut}, applied to $B$, the vertices $p$ and $q$
cannot be separated from $x$ in $B$ by at most two vertices. Since every
separator of $J$ of size at most three contains $w$, they are good in $J$.
\end{itemize}
Goodness stays true under later replacements, by Lemma~\ref{lem:retain}.

Now consider a bad component after such a replacement. It cannot contain any of
the good boundary vertices. Therefore it contains no end of a newly added
clique edge, and it has no old edge into the deleted part, whose neighbors are
$w,p,q$. So it was already \emph{the same component, with the same vertex
set}, before the replacement.

Consequently it keeps its old marker. This proves, by induction, that every
bad component during restoration contains an original marker that is still
present. At the start of restoration this holds by the last statement of
Section~\ref{sec:guarantees}, since by properties 1 and 2 every bad component
comes from a separator $\{w,p,q\}$. After the marker list is exhausted:
\begin{itemize}[nosep]
\item a remaining marker whose flow succeeded is good, because a separator of
$J$ of size at most three would contain $w$, and the rest of it would
separate $z$ from $x$ in $B$;
\item a marker whose flow failed was deleted;
\item a skipped cycle marker is good, because it is adjacent to $x$.
\end{itemize}
A remaining bad component would therefore contain a good vertex, which is a
contradiction. So no bad component remains, and the instance is clean.

\begin{lemma}\label{lem:restore}
Restoration returns a clean instance that has a crossing if and only if the
contracted instance has one, and it uses at most $|M|$ flow queries for three
paths.
\end{lemma}

The bound holds even for a trial that is rejected in the end. There is no
hidden search over all remaining vertices.

\section{The complete contraction algorithm}\label{sec:loop}

After the initial cleaning, call $\Classify$.
\begin{itemize}[nosep]
\item On \ans{NO}, report that $C$ is crossless.
\item On \ans{WITNESS}, lift the crossing through the recorded initial
replacements and return it.
\item On \ans{YES}, the current graph is clean, nonplanar and crossed.
\end{itemize}

We repeat the following iteration while in the last state.

\paragraph{One iteration.}
\begin{enumerate}[nosep]
\item \emph{Sparsify} as in Section~\ref{sec:sparsify}, if needed. The graph
stays clean, nonplanar and crossed, and now $M\le16N$.
\item Choose any five distinct vertices of degree at most $64$.
\item For each chosen vertex $u$:
\begin{itemize}[nosep]
\item if $u\notin V(D)$, try every incident edge $uv$;
\item if $u\in V(D)$, try its two cycle edges, except that contractions of
cycle edges are skipped when $|D|=4$.
\end{itemize}
\item For each trial:
\begin{itemize}[nosep]
\item copy the current instance;
\item contract $uv$;
\item restore cleanliness as in Section~\ref{sec:repair}, using the markers
$N_K(u)\setminus\{v\}$;
\item call $\Classify$ on the restored instance.
\end{itemize}
\item On \ans{NO}, discard that trial.
\item On \ans{WITNESS}, lift the crossing through the replacements and the
contraction of that trial, then through all earlier accepted operations, and
return it.
\item On \ans{YES}, accept that trial and begin the next iteration.
\end{enumerate}

All choices in the graph and all orders of adjacency lists may be fixed
arbitrarily.

\subsection{Five vertices of low degree always exist}\label{sec:five}

A nonplanar graph has at least five vertices. For $5\le N\le9$, every vertex has
degree at most eight. For $N\ge10$, the number of vertices of degree at least
$65$ is at most
\begin{equation*}
\frac{2M}{65}\le\frac{32N}{65}<\frac N2.
\end{equation*}
Thus more than $N/2\ge5$ vertices have degree at most $64$. No exhaustive base
case for small instances is needed.

\subsection{At least one trial is positive}\label{sec:positive}

Fix any crossing of $D$. It has exactly four ends. So at least one of the five
chosen vertices, say $u$, is not one of these ends. There are three cases.
\begin{itemize}[nosep]
\item \emph{The vertex $u$ is off the cycle and lies on a crossing path.}
Contract an edge of that path at $u$. The crossing survives.
\item \emph{The vertex $u$ is off the cycle and lies on neither path.}
Contract any edge at $u$. The crossing survives. Such an edge exists, because
cleanliness implies that every vertex off the cycle has degree at least four.
\item \emph{The vertex $u$ is on the cycle.} Then neither crossing path uses
it. Contract either cycle edge at $u$. The cyclic order of the four ends of the
crossing does not change. Here $|D|\ge5$, since $D$ contains the four ends and
$u$, so this trial is not skipped.
\end{itemize}
Restoration preserves the existence of a crossing, by
Lemma~\ref{lem:replace}. Since the classification is exact, it accepts at least
one trial.

Every accepted trial contracts an edge and so decreases the number of vertices.
Thus the procedure terminates. It cannot end in a nonplanar positive state with
fewer than five vertices, so before that happens the classification returns a
crossing.

\section{Lifting contractions}\label{sec:lift}

For completeness, here are all the lifts through contractions.

The record stores the graph and the cycle before the contraction, and the edge
$uv$. At most one returned crossing path uses the new vertex $w$, since the
two paths are disjoint.

\paragraph{Both ends off the cycle.}
If $w$ is unused, keep the paths. Otherwise let $a$ and $b$ be the vertices
before and after $w$ on its path. In the recorded graph, choose
\begin{equation*}
u_a\in\{u,v\}\text{ adjacent to }a,
\qquad
u_b\in\{u,v\}\text{ adjacent to }b.
\end{equation*}
If $u_a=u_b$, replace $w$ by that vertex. Otherwise replace it by $u_a,u_b$,
using the edge $uv$. The added vertices do not occur on the other path or
elsewhere on this path.

\paragraph{One end on the cycle.}
The algorithm orients this case so that $u\notin V(D)$ and $v\in V(D)$, and
keeps $w=v$. If $w$ is used, it is an end of a crossing path. Suppose that the
path starts at $w$. If its first edge comes from $v$ in the recorded graph,
keep the end $v$. If it comes from $u$, replace the beginning by $v,u$,
followed by that edge. For a path that ends at $w$, reverse the order. The only
vertex added inside the path is $u$, which is off the cycle.

\paragraph{Both ends on the cycle.}
Then $u$ and $v$ were consecutive on the recorded cycle. If $w$ is used, it is
an end of a crossing path. Choose whichever of $u,v$ is adjacent, in the
recorded graph, to the next vertex of the path, and use it as the end.
\emph{Do not insert the cycle edge $uv$ into the path.} Replacing one cycle
position by one of two consecutive positions keeps the cyclic order of the
four distinct ends, since no other end lies between those two positions.

Thus every contraction lifts without a search. Reversing all accepted
operations gives two disjoint $C$-paths in the original graph, with
alternating ends.

\section{Search count and total running time}\label{sec:time}

Let $n_0$ and $m_0$ be the numbers of vertices and edges of the component that
contains the original cycle. This component is connected and contains a
cycle, so $m_0\ge n_0$.

\begin{center}
\begin{tabular}{lr}
\toprule
Operation & Number\\
\midrule
Initial flow queries for four paths & at most $n_0$\\
Accepted contraction iterations & at most $n_0$\\
Trials per iteration & at most $5\cdot64=320$\\
Flow queries for three paths per trial & at most $63$\\
BFS forests per sparsification & four\\
Boolean planarity tests & at most $320n_0+1$\\
Searches constructing the final bridge crossing & two\\
Searches reversing separator replacements & at most $n_0$\\
\bottomrule
\end{tabular}
\end{center}

A flow query for three paths uses at most three residual searches. On failure,
finding the maximum cut side and the relevant components takes a bounded
number of further linear traversals. Thus the total number of residual and
ordinary graph searches is $O(n_0)$.

The planarity tests are \emph{separate linear-time calls}. We do not count them
as ordinary path searches.

Every accepted separator replacement deletes at least one vertex and adds at
most three edges. Contractions add no edges before simplification. Temporary
trials make only a bounded number of replacements. Consequently every graph
that is searched has size
\begin{equation*}
O(n_0+m_0).
\end{equation*}

All remaining work is linear per bounded-flow query or per trial: graph
copies, contraction and simplification, bridge construction, counting sorts,
stack operations, certificate construction, and path splicing. Storing
complete snapshots is allowed, since the stack of accepted operations contains
only $O(n_0)$ operations.

Therefore the total time is
\begin{equation*}
O\bigl(n_0(n_0+m_0)\bigr)
=
O(n_0m_0)
\subseteq O(nm).
\end{equation*}

\section{Assumptions, critical choices, and tests}\label{sec:remarks}

The algorithm runs on the given pair $(G,C)$. It assumes only that $G$ is finite
and simple and that $C$ is a cycle of $G$.

The following choices matter:
\begin{itemize}[nosep]
\item BFS (scan-first) forests for the certificate. Arbitrary spanning forests
do not suffice.
\item A minimum two-cut with \emph{maximum open source side} during
restoration.
\item Contractions only of edges at a vertex off the cycle, or of edges between
consecutive cycle vertices.
\end{itemize}

The file \texttt{crossing\_algorithm.py} in this folder is a reference
implementation in Python. It has the interface
\begin{verbatim}
from crossing_algorithm import find_crossing

result = find_crossing(n, edges, cycle)
# None, or (path1, path2)
\end{verbatim}
It uses the planarity test of the companion note \cite{HJ26}, from the file
\texttt{planarity\_linear.py} in the folder \texttt{docs/planarity}, and no
other library outside the Python standard library. The default call follows
the algorithm above. After sparsifying, it also calls $\BridgeCross$ once more,
which is not in the algorithm above and does not change the bounds. The
optional mode \texttt{audit=True} checks cleanliness again after every trial.
This mode is for testing, and it does not meet the bound $O(nm)$. Python sets
and dictionaries are not the arrays of a word RAM, so the implementation
illustrates the algorithm but does not prove its time bound.

We tested the implementation as follows. On 9{,}000 random instances with at
most $19$ vertices, the answer agreed with an exhaustive search for a crossing,
with the audit mode on. These instances include 3{,}731 runs of the
contraction loop with 15{,}296 trials, many of them with restoration
replacements. On 40 dense random instances with $36$ to $45$ vertices, the
sparsification step ran 35 times. Every returned crossing is checked against
the original graph before it is returned. Separately, on 100{,}000 random
instances, $\BridgeCross$ agreed with a direct check for alternating
attachments of two distinct bridges. These tests check the implementation. The
correctness and the time bound follow from the lemmas above.

\section{The Lean formalization}\label{sec:lean}

The Lean formalization of this note is in the repository \texttt{GraphMinor}.
\begin{itemize}[nosep]
\item The folder \path{GraphMinor/Crossing} proves the mathematics of
Sections~\ref{sec:facts} to~\ref{sec:lift}, apart from Fact~P. The file
\path{GraphMinor/Crossing/Statements.lean} states the claims of
Sections~\ref{sec:root}, \ref{sec:replace}, \ref{sec:clean}, \ref{sec:sparse},
\ref{sec:twocut}, \ref{sec:repair}, \ref{sec:loop} and~\ref{sec:lift} as
propositions, each with the passage it formalizes. Fact~T is proved in
\path{GraphMinor/Crossing/PlanarFactT.lean} from the formalization of
Humeau and Pous in the folder \path{GraphMinor/HP25}.
\item In the folder \path{GraphMinor/Compute}, the files
\path{Crossing*.lean} and \path{SpecsCrossing*.lean} give word-RAM programs
for all steps and prove their correctness and running times. The whole
algorithm is \texttt{CrossingExtractorStatement}, proved as
\texttt{crossingExtractor\_of}.
\item The file \path{GraphMinor/Compute/SpecsCrossingReductionProved.lean}
derives from it the two paths statement \texttt{Shiloach80Statement} used in
(8.1), by the reduction recalled in the introduction. The file
\path{GraphMinor/Compute/DisjointPathsFinal.lean} proves the main result of
Graph Minors XIII, \texttt{disjointPathsStatement\_proved}, with no hypotheses.
\end{itemize}

The Lean formalization differs from this note in the following places.
\begin{itemize}[nosep]
\item \emph{The model of computation.} The note counts searches and argues
that the rest of the work is linear per search. The Lean proves executions of
fixed word-RAM programs, with explicit word sizes, memory layouts and record
buffers, and bounds their numbers of steps. A search count alone does not show
that a program with bounded words meets the time bound.
\item \emph{The planarity test.} Fact~P is the Lean statement
\texttt{PlanarityTestStatement}. It is not assumed: it is proved, as
\texttt{planarityTestStatement\_proved}, from the companion note \cite{HJ26}.
\item \emph{The use in Graph Minors XIII.} The algorithm of this note returns
paths. The Lean uses it for the step of (8.1), and then assembles the
algorithm for the main result of Graph Minors XIII, which decides the disjoint
paths problem for a fixed number of terminal pairs. That final program returns
a Boolean answer, as the problem is stated in Graph Minors XIII, page~65. It is
not required to return the paths.
\end{itemize}

{\small
\begin{thebibliography}{GM13}

\bibitem[CKT93]{CKT93}
J.~Cheriyan, M.-Y.~Kao and R.~Thurimella,
Scan-first search and sparse certificates: an improved parallel algorithm for
$k$-vertex connectivity,
\emph{SIAM J. Comput.} \textbf{22} (1993), 157--174.

\bibitem[GM13]{GM13}
N.~Robertson and P.~D.~Seymour,
Graph Minors. XIII. The disjoint paths problem,
\emph{J. Combin. Theory Ser. B} \textbf{63} (1995), 65--110.

\bibitem[HJ26]{HJ26}
M.~Hajiaghayi and M.~JafariRaviz,
A Boolean linear-time planarity test,
note, November 2026, file \path{docs/planarity/planarity_test.tex} of the
repository \texttt{GraphMinor}.

\bibitem[HP25]{HP25}
S.~Humeau and D.~Pous,
On the two paths theorem and the two disjoint paths problem,
arXiv:2505.16431v2 (14 August 2025).

\bibitem[HT74]{HT74}
J.~Hopcroft and R.~Tarjan,
Efficient planarity testing,
\emph{J. ACM} \textbf{21} (1974), 549--568.

\bibitem[Sh80]{Sh80}
Y.~Shiloach,
A polynomial solution to the undirected two paths problem,
\emph{J. ACM} \textbf{27} (1980), 445--456.

\end{thebibliography}}

\end{document}
